\documentclass[Total.tex]{subfiles}

\begin{document}
	\chapter{Riemann Surfaces}
		\section{Definitions}
		
			\begin{Def}
				A \textbf{Riemann surface} is one-dimensional complex=\underline{conneected} $1$-dimensional complex analytic manifold.
				
				That is, equipped with maximal charts $\{(U_\alpha,z_\alpha)\}_{\alpha\in A}$ on $M$ 
				
				\begin{gather*}
					z_\alpha:U_\alpha\rightarrow \mathbb{C}
				\end{gather*}
				
				Then transition functions are holomorphic
				
				\begin{gather*}
					f_{\alpha\beta}=z_\alpha\circ z_\beta^{-1}:;U_\alpha\cap U_\beta=\varnothing
				\end{gather*}
				
				The charts are called \textbf{transition functions}.
			\end{Def}
			
			Notice that: The set of holomorphic functions forms a pseudogroup under composition. Taking
			
			\begin{gather*}
				\mathcal{H}(X)=\{f:U\rightarrow V|f\mbox{ is biholomorphism};U,V\subseteq X\mbox{ are open}\}
			\end{gather*}
			
			Then,
			
			\begin{itemize}
				\item 
				
				\item 
				
				\item 
				
				\item 
			\end{itemize}
			
			\begin{Def}
				Classically, a compact Riemann surface is called \textbf{closed}, elsewise is \textbf{open}.
			\end{Def}
			
			\begin{lemma}
				内容...
			\end{lemma}
			
			\begin{proof}
				内容...
			\end{proof}
			
			
			\subsection{Holomorphic Mapping}
			
		\section{Topology of Riemann Surface}
			\subsection{Fundamental Groups}
			
			\subsection{Covering Manifolds}
			
			\subsection{Riemann-Hurwitz Relation}
			
				\begin{theorem}
					We have
					
					\begin{gather*}
						g=n(\gamma-1)+1+B/2
					\end{gather*}
				\end{theorem}
		
		\section{Differential Forms}
		
		\section{Integration Forms}
	\chapter{Analytic Continuation, Covering Surfaces and Algebraic Functions}
		\section{Analytic Continuation}
			\begin{Def}
				\begin{itemize}
					\item (Function Element) Let $M,N$ be fixed Riemann surfaces. A \textbf{function element} is pair
					\begin{gather*}
						(f,D)
					\end{gather*}
					
					where:
					\begin{itemize}
						\item $D$ is a connected, open nonempty subset of $M$;
						\item $f:D\rightarrow N$ is analytic. 
					\end{itemize}
					
					\item Two elements $(f_1,D_1),(f_2,D_2)$ are \textbf{direct analytic continuation} of each other if
					\begin{gather*}
						D_1\cap D_2\not=\emptyset\qquad f_1|D_1\cap D_2=f_2|D_1\cap D_2
					\end{gather*}
					
					Use the following notion $(f_1,D_1)\simeq (f_2,D_2)$.
				\end{itemize}
			\end{Def}
			
			\begin{Def}
				(Analytic Continuations)
				\begin{itemize}
					\item We say $(f,D)$, $(g,\tilde{D})$ are \textbf{analytic continuation} if there exists function elements $(f_j,D_j);1\leq j\leq J$: 
					\begin{gather*}
						(f_1,D_1)=(f,D)\qquad (f_J,D_J)=(g,\tilde{J})\qquad (f_j,D_j)\simeq (f_{j+1},D_{j+1})
					\end{gather*}
					\item A \textbf{compelete analytic continuation} is an equivalence class of suc $(f,D)$. By analytic continuation.
				\end{itemize}
			\end{Def}
			
			\begin{lemma}
				(Analytic Continuation$\Rightarrow$ Path) Suppose $(f,D),(g,\tilde{D})$ is an analytic continuation. Then, there exists a path
				\begin{gather*}
					\gamma:[0,1]\rightarrow M
				\end{gather*}
				
				Connecting the points of $D$, with following property: Exists
				\begin{gather*}
					0<t_1<\dots<t_m=1
				\end{gather*}
				
				where: Exists $B(\gamma(t_i),r_i)$
			\end{lemma}
			\begin{proof}
				内容...
			\end{proof}
			
			\begin{corollary}
				内容...
			\end{corollary}
			\begin{proof}
				内容...
			\end{proof}
			
			\section{The Unramified Riemann Surface of An Analytic Germ}
			\section{The Ramified Riemann Surface of An Analytic Germ}
			\section{Algebraic Germs and Functions}
			\section{title}
	\chapter{Stein Manifold}
	\chapter{Compact Riemann Surface}
		\section{Existence Theorem}
\end{document}